\begin{table}[ht]
\centering
\small
\caption{Performance under observation noise and unreliable action execution. The first row of each block is the control condition, in which both manipulations are switched off and the published run is recovered.}
\label{tab:noise}
\begin{tabular}{llrrrrrr}
\toprule
Agent & Condition & Seeds & 0--5 min & 6--20 min & Median & $\Delta$ & $p$ \\
\midrule
CFL-4 & control & 100 & 1.08 & 2.39 & 1.83 & +1.31 & $<$0.001 \\
CFL-4 & obs. $\varepsilon = 0.05$ & 100 & 2.40 & 6.29 & 2.81 & +3.89 & $<$0.001 \\
CFL-4 & obs. $\varepsilon = 0.15$ & 100 & 1.94 & 4.75 & 3.27 & +2.81 & $<$0.001 \\
CFL-4 & act. $\eta = 0.10$ & 100 & 1.01 & 1.46 & 1.24 & +0.44 & $<$0.001 \\
DP-5 & control & 6 & 0.91 & 1.30 & 1.21 & +0.38 & 0.273 \\
DP-5 & obs. $\varepsilon = 0.05$ & 6 & 0.85 & 1.12 & 1.10 & +0.27 & 0.117 \\
DP-5 & obs. $\varepsilon = 0.15$ & 6 & 0.85 & 0.90 & 0.94 & +0.04 & 0.812 \\
DP-5 & act. $\eta = 0.10$ & 6 & 1.23 & 0.73 & 0.73 & -0.49 & 0.018 \\
\bottomrule
\end{tabular}
\\[4pt]\footnotesize{The median is given because the distribution of rally length over seeds is heavy tailed for CFL-4, so the mean alone would overstate the typical run. The DP-5 rows rest on six seeds and are indicative only; no difference between its observation-noise conditions is distinguishable from zero.}
\end{table}
